% ============================================================================
% Section 2: Mathematical Preliminaries
% ============================================================================
\section{Mathematical Preliminaries}
\label{sec:math_prelim}

\subsection{The Dedekind Eta Function}

\begin{definition}[Dedekind Eta Function~\cite{dedekind_1877}]
\label{def:dedekind_eta}
For $\tau \in \mathfrak{H} = \{\tau \in \CC : \mathrm{Im}(\tau) > 0\}$
and $q = e^{2\pi i \tau}$, the \emph{Dedekind eta function} is
\begin{equation}
  \eta(\tau) = q^{1/24} \prod_{n=1}^{\infty} (1 - q^n).
\end{equation}
It is a modular form of weight $1/2$ for the full modular group
$\mathrm{SL}_2(\ZZ)$ with a multiplier system~\cite{ono2004}.
\end{definition}

\begin{definition}[Eta-Quotient~\cite{martin_1996,gordon_hughes_1993}]
\label{def:eta_quotient}
For a positive integer $N$ (the \emph{level}) and integers
$r_d \in \ZZ$ indexed by divisors $d \mid N$, an
\emph{eta-quotient} of level $N$ is the function
\begin{equation}
  \label{eq:eta_quotient}
  f(\tau) = \prod_{d \mid N} \eta(d\tau)^{r_d}.
\end{equation}
The vector $(r_d)_{d \mid N}$ is called the
\emph{exponent vector}.
\end{definition}

\subsection{Modularity Conditions}

An eta-quotient $f(\tau) = \prod_{d \mid N} \eta(d\tau)^{r_d}$
is a modular form of weight $k = \frac{1}{2}\sum_{d \mid N} r_d$
for $\Gamma_0(N)$ if and only if the following conditions are
satisfied~\cite{ono2004,martin_1996}:
\begin{enumerate}[label=(M\arabic*)]
  \item $k = \frac{1}{2}\sum_{d \mid N} r_d \in \ZZ$
        \hfill (integer weight),
  \item $\sum_{d \mid N} d \cdot r_d \equiv 0 \pmod{24}$
        \hfill (cusp at $\infty$),
  \item $\sum_{d \mid N} \frac{N}{d} \cdot r_d \equiv 0 \pmod{24}$
        \hfill (cusp at $0$),
  \item For each cusp $s$ of $\Gamma_0(N)$:
        $\sum_{d \mid N} \frac{\gcd(d,s)^2}{d} \cdot r_d \geq 0$
        \hfill (holomorphicity).
\end{enumerate}

\subsection{Physical Invariants of an Eta-Quotient}

\begin{definition}[Effective Central Charge]
\label{def:c_eff}
For an eta-quotient with exponent vector $(r_d)_{d \mid N}$, the
\emph{effective central charge} is
\begin{equation}
  \label{eq:c_eff}
  \ceff = \sum_{d \mid N} \frac{r_d}{d}.
\end{equation}
In the context of two-dimensional conformal field theory (CFT),
$\ceff$ governs the asymptotic density of states via the Cardy
formula~\cite{cardy_1986}.
\end{definition}

\begin{definition}[Modular Weight]
\label{def:weight}
The \emph{modular weight} of the eta-quotient is
\begin{equation}
  k = \frac{1}{2} \sum_{d \mid N} r_d.
\end{equation}
\end{definition}

\begin{definition}[Leading Power]
\label{def:leading_power}
The \emph{leading power} (order of vanishing at $q = 0$) is
\begin{equation}
  \label{eq:leading_power}
  p = \frac{1}{24} \sum_{d \mid N} d \cdot r_d.
\end{equation}
\end{definition}

\subsection{The Cardy Formula and BPS State Entropy}

In a unitary 2D CFT with central charge $c$ and effective
central charge $\ceff$, the asymptotic number of states at
excitation level $n$ is given by the Cardy
formula~\cite{cardy_1986}:
\begin{equation}
  \label{eq:cardy}
  \rho(n) \sim \exp\left(2\pi \sqrt{\frac{\ceff \cdot n}{6}}\right),
  \qquad n \to \infty.
\end{equation}
The BPS state entropy in the Strominger--Vafa
framework~\cite{strominger_vafa_1996} is
\begin{equation}
  \label{eq:bps_entropy}
  S_{\mathrm{BPS}} = 2\pi \sqrt{\frac{\ceff}{6}} \cdot \sqrt{n}
  = 2\pi \sqrt{\frac{c_{\mathrm{eff}} \cdot n}{6}}.
\end{equation}

\subsection{Mock Theta Functions}

\begin{definition}[Mock Theta Function~\cite{zwegers2002,bringmann_ono_2006}]
In Ramanujan's final letter to Hardy (1920), he introduced
functions that ``mock'' the behavior of modular forms near
the boundary of the unit disk.  Zwegers~\cite{zwegers2002}
and Bringmann--Ono~\cite{bringmann_ono_2006} established that
mock theta functions are the holomorphic parts of
weight-$1/2$ \emph{harmonic Maass forms}.  The
\emph{shadow} of a mock theta function is a unary theta
series or eta-quotient that encodes the non-holomorphic
completion.
\end{definition}

\begin{remark}
The connection between mock theta functions and
eta-quotients is precisely the bridge exploited by the
\RAMA{} engine: the evolutionary search targets
eta-quotient shadows of mock modular forms extracted
from Ramanujan's manuscript pages.
\end{remark}

\subsection{The Hardy--Ramanujan--Rademacher Expansion}

The partition function $p(n)$ satisfies
the Hardy--Ramanujan asymptotic
formula~\cite{hardy_ramanujan_1918}:
\begin{equation}
  p(n) \sim \frac{1}{4n\sqrt{3}}
  \exp\left(\pi\sqrt{\frac{2n}{3}}\right),
  \qquad n \to \infty.
\end{equation}
Rademacher~\cite{rademacher_1937} refined this to an exact
convergent series.  The coefficients of general
eta-quotients admit analogous exact expansions, which
provide the theoretical foundation for comparing
computationally discovered $q$-series coefficients against
analytic predictions.
